\section{Формальный финал}

\begin{frame}
    \frametitle{Положения, выносимые на защиту}
    \small
    \begin{enumerate}
        \item Разработана методика построения, исследования и аппроксимации модельного многообразия решений энергетически оптимальных задач встречи с идеально регулируемой тягой.
        \item Получены конечные формулы, задающие аппроксимированное многообразие решений в плоской круговой модели движения небесных тел.
        \item Получены конечные формулы и линейные преобразования начальных значений сопряжённых переменных, связывающие параметры пространственной и эллиптической конечной орбиты с начальными приближениями расширенной системы.
        \item Разработан и апробирован алгоритм выбора начального приближения для энергетически оптимальных межпланетных перелётов с непрерывной тягой.
    \end{enumerate}
\end{frame}

\begin{frame}
    \frametitle{Публикации: ВАК / Scopus / РИД}
    \begin{columns}[T,onlytextwidth]
        \begin{column}{0.5\textwidth}
            \begin{itemize}
                \item По теме диссертации опубликовано 19 научных работ.
                \item 6 работ опубликованы в изданиях, рекомендованных ВАК РФ.
                \item 2 статьи индексируются в международных базах Scopus и Web of Science.
            \end{itemize}
        \end{column}
        \begin{column}{0.46\textwidth}
            \begin{block}{Программная реализация}
                Получено свидетельство о государственной регистрации программы для ЭВМ, относящейся к реализации разработанного метода.
            \end{block}
        \end{column}
    \end{columns}
\end{frame}

\begin{frame}
    \frametitle{Публикации: РИНЦ и материалы конференций}
    \begin{itemize}
        \item 13 работ опубликованы в изданиях и материалах научных конференций, индексируемых в РИНЦ.
        \item Публикации покрывают основные блоки работы: модельное многообразие, конечные формулы, \(R\)-преобразования, прикладные перелёты Земля--Марс и Земля--Венера.
        \item Полные библиографические описания приведены в тексте диссертации и списке работ автора.
    \end{itemize}
\end{frame}

\begin{frame}
    \frametitle{Апробация}
    \begin{columns}[T,onlytextwidth]
        \begin{column}{0.52\textwidth}
            \textbf{Конференции}
            \begin{itemize}
                \item МФТИ, 2020, 2021, 2023, 2025;
                \item Академические чтения по космонавтике, 2022, 2024, 2025;
                \item «Авиация и космонавтика», 2024;
                \item «Новые горизонты прикладной математики», 2024, 2025.
            \end{itemize}
        \end{column}
        \begin{column}{0.44\textwidth}
            \textbf{Семинары}
            \begin{itemize}
                \item ИПМ им. М.\,В. Келдыша РАН;
                \item ИКИ РАН;
                \item ИПМех РАН;
                \item МАИ.
            \end{itemize}
        \end{column}
    \end{columns}
\end{frame}

\begin{frame}[plain]
    \begin{center}
        \Huge
        Спасибо за внимание!
    \end{center}
\end{frame}
